\section{结论}

本文围绕轴承预测性维护，完成了 XJTU-SY 和 PHM2012 两个数据集上的三任务实验闭环。通过统一的数据加载、按轴承划分、manual\_basic 特征分析、\code{linear_rul_norm} 标签构造和 GRU 序列基线训练，本文将剩余寿命预测、健康阶段识别和早期故障检测纳入同一套可复查流程。

实验结果表明，GRU 序列基线能够在 EarlyFault 任务上取得较稳定表现，并在 RUL 任务上形成可解释的寿命趋势预测。与此同时，预测质量排查也提醒我们：RUL 泛化仍是难点，尤其在跨轴承、跨工况情况下不能夸大精度。本文因此将主线结果、对照模型实验、论文复现实验和演示视频分开表述，保证结论边界清楚。

总体而言，本项目完成了从原始振动信号到特征分析、标签构造、模型训练、结果审计和中文演示材料的完整链路。结合论文复现实验和训练过程视频，当前材料已经具备结题报告、答辩 PPT 和现场演示的基础。
